বাক্যের সেটগুলি এক অর্থে সেই গঠনগুলিকে বর্ণনা করে, যেগুলিতে
বাক্যগুলি একসঙ্গে সত্য। ওই গঠনগুলিই তাদের \emph{মডেল}।
উল্টো দিক থেকে, একটি গঠন বা গঠনের সেট দিয়ে শুরু করলে, তারা
যে বাক্যগুলির মডেল সেই বাক্যগুলির সেটটিও বিবেচনা করা যায়।
সেটিই ওই গঠন বা গঠনের সেটের \emph{তত্ত্ব}। এমন যেকোনো বাক্যের
সেটে তা থেকে অনুসৃত প্রত্যেক বাক্য আগে থেকেই থাকে। অর্থাৎ,
সেটগুলি \emph{বদ্ধ}। আরও সাধারণভাবে, সেট $\Gamma$-কে তত্ত্ব
বলি যদি সেটি অর্থগত অনুসরণের অধীনে বদ্ধ হয়। আর $\Gamma$-কে
~$\Delta$ দ্বারা \emph{স্বতঃসিদ্ধায়িত} বলি যদি~$\Delta$ থেকে
অনুসৃত সব বাক্য নিয়ে~$\Gamma$ গঠিত হয়।

গণিতে তত্ত্বের অনেক উদাহরণ আছে। যেমন, রৈখিক ক্রমের তত্ত্ব,
গোষ্ঠীর তত্ত্ব, অথবা পাটিগণিতের নানা তত্ত্ব, যেমন পেয়ানোর
স্বতঃসিদ্ধগুলি দ্বারা স্বতঃসিদ্ধায়িত তত্ত্ব। অন্যান্য
বিদ্যাশাখাতেও গুরুত্বপূর্ণ তত্ত্বের অনেক উদাহরণ আছে। যেমন,
সম্পর্কভিত্তিক ডেটাবেসকে তত্ত্ব হিসেবে ভাবা যায়, এবং
অধিবিদ্যায় অংশ-সম্পর্কের এমন তত্ত্ব আলোচিত হয়, যেগুলি
স্বতঃসিদ্ধায়িত করা যায়।

কোনো তত্ত্ব স্থির করে অধ্যয়ন করতে গেলে একটি গুরুত্বপূর্ণ
প্রশ্ন হলো, আমরা তত্ত্বটিতে যে সব কথা বলতে চাই, সেগুলি প্রকাশ করার
মতো প্রকাশক্ষমতা তার ভাষার আছে কি না। আরেকটি প্রশ্ন হলো,
আমরা যা যা প্রমাণ করতে চাই, তা প্রমাণ করার মতো শক্তি তত্ত্বটির
আছে কি না। কোনো সম্পর্ক \emph{প্রকাশ} করতে হলে প্রয়োজনীয়
সংখ্যক মুক্ত চলরাশি-সহ একটি সূত্র লাগে। \emph{সেটতত্ত্বে}
সম্পর্ক সংকেত হিসেবে কেবল $\in$ আছে, তবু $x \subseteq y$ প্রকাশ
করা যায় $\lforall[u][(u \in x \lif u \in y)]$ সূত্রটি দিয়ে। বস্তুত
\emph{জার্মেলো--ফ্রেঙ্কেল সেটতত্ত্ব}~$\Log{ZFC}$ এতটাই শক্তিশালী যে
অল্প কয়েকটি স্বতঃসিদ্ধ এবং ক্রমশ জটিলতর হয়ে ওঠা সংজ্ঞার
(যেমন~$\subseteq$-এর সংজ্ঞা) একটি শৃঙ্খল ব্যবহার করে (প্রায়)
প্রতিটি গাণিতিক দাবি প্রকাশ এবং (প্রায়) প্রতিটি গাণিতিক উপপাদ্য
প্রমাণ করা যায়।
